\section{Limitations}
\label{sec:limitations}

\textbf{Data size and label noise.} The benchmark has 53 patients; the test split holds 206
criteria from 8 patients, no \emph{excluded} example, and only 3--7 examples of several other
classes, so per-class test figures have wide uncertainty and the test set is dominated by a few
patients and trials (one patient contributes 14 of the 25 largest-group errors). The error analysis
shows that annotators applied the \emph{not excluded} versus \emph{not enough information}
distinction inconsistently for conditions the note does not mention; part of the measured error
reflects this ambiguity rather than model failure, and conversely the model has learned a convention
(absence $\Rightarrow$ not excluded) that a screening tool should not apply blindly.

\textbf{Model capacity and reasoning.} A 1.5B model does not perform the clinical inferences and
counting that some criteria require (Table~\ref{tab:errors}), and it essentially never predicts the
rarest classes. Explanations are assembled from cited sentences rather than generated, which avoids
hallucinated rationales but cannot express the reasoning linking evidence to label. The model judges
each criterion independently and ignores compound logic across criteria.

\textbf{Selection and evaluation.} Model selection used a validation set too small to separate the
variants statistically; the chosen configuration's advantage is not significant, although seed
repeats show it is stable. We evaluated a single dataset; generalization to real clinical notes,
which are longer, noisier and not pre-split into sentences, is untested. The application's sentence
splitter for unseen notes is rule-based and may split or merge sentences differently from the
dataset's annotation.

\textbf{Inference implementations.} Our decoders run at batch size 1 and do not implement paged or
quantized caches. The tiled attention is an educational Python implementation that saves memory but
is much slower than fused kernels. Speculative decoding gives no speed-up in our setting because the
draft is not cheap enough relative to the target at batch size 1; we did not evaluate larger target
models, where it would be expected to help. Timings were measured on one T4 and one laptop CPU and
include Kaggle's shared-infrastructure variability.

\textbf{Clinical use.} The system is an educational screening prototype. It has not been validated
clinically, is not a medical device, and must not be used for enrollment decisions; every result
requires review by qualified staff.
